\hypertarget{appendix-05f39b10364d}{}
\section*{Minimum signed value}
\section{Minimum signed value}\label{minimum-signed-value}

\textbf{Curriculum depth:} Introduced · \textbf{First seen:} C002 · \textbf{Latest development:} C002

\subsection{Definition and mechanism}\label{definition-and-mechanism}

Two's-complement range has one extra negative value, so \texttt{INT\_MIN} cannot be negated or represented by same-type absolute value.

N-bit two's complement spans −2ᴺ⁻¹ through 2ᴺ⁻¹−1. Eight-bit minimum is −128, whose positive counterpart 128 does not fit. Negating int minimum or dividing it by −1 overflows the operation type. Narrow int8\_t operands may promote first, so the original storage width alone does not establish UB.

\subsection{Prerequisites}\label{prerequisites}

\begin{itemize}
\tightlist
\item
  \hyperref[appendix-4fbc03d20829]{Two\textquotesingle s complement} --- For raw unsigned value u in N bits, signed decoding is u when u\textless2ᴺ⁻¹, otherwise u−2ᴺ.
\end{itemize}

\subsection{Used by}\label{used-by}

\begin{itemize}
\item
  \hyperref[appendix-3333db47830f]{Fixed-point rescaling} --- uses this mechanism to define its representation, error or validation contract.
\item
  \hyperref[appendix-1a599c924a4e]{Signed integer overflow} --- uses minimum signed value as part of its representation or reasoning.
\end{itemize}

\subsection{Recall and next depth}\label{recall-and-next-depth}

Explain minimum signed value without the label, work the example above, then state the assumption whose removal would change the result.

\textbf{Pending:} Apply and reconstruct this concept independently.

\subsection{Source and navigation}\label{source-and-navigation}

\href{Continuous\%20Notes/Computing\%20Foundations\%20-\%20From\%20Bits\%20to\%20C++\%20and\%20Systems\%20APIs\#Arithmetic\%20becomes\%20a\%20policy}{C002 teaching chapter} · \hyperref[appendix-fd090d982a15]{Part definitions and mechanisms} · \href{}{Dependency Map (02-Foundations-II, appendix-4ad47bdd269a)}
